\section{Tehničke metode maksimizacije korisničkog angažovanja}

Iza prividno jednostavnog i intuitivnog korisničkog interfejsa, savremene platforme koriste složene algoritamske sisteme za maksimizaciju korisničkog angažovanja. Ne maksimizuje se samo vreme koje korisnik provodi pred ekranom, već i čitav niz interakcija sa platformom. O korisnicima se prikupljaju podaci koji zatim imaju dvojaku funkciju. S jedne strane, koriste se za profilisanje korisnika, ciljano oglašavanje i poslovne odluke zasnovane na analitici ponašanja. S druge strane, isti podaci se koriste i za dalje prilagođavanje algoritama preporuke, čime se stvara povratna sprega: korisnik reaguje na sadržaj, sistem tu reakciju beleži, a zatim naredni sadržaj prilagođava tako da poveća verovatnoću nove reakcije.

Zbog toga se sistemi preporuke ne mogu posmatrati samo kao neutralni alati za sortiranje informacija. Oni su tehnički mehanizmi kroz koje platforma aktivno oblikuje korisničko iskustvo. Pojedini istraživači opisuju ovaj proces kao deo šire ekonomije pažnje: savremeni poslovni modeli digitalnih platformi koriste veštačku inteligenciju i analitiku velikih podataka kako bi povećali korisničko angažovanje, pri čemu se stvara ciklus privlačenja pažnje i ekstrakcije podataka \cite{gonzalez2024attention}. Drugim rečima, korisnik svojim ponašanjem stalno proizvodi nove podatke, a ti podaci zatim služe da se naredni sadržaj još preciznije prilagodi njegovim reakcijama.

U ovom poglavlju biće obrađene tri tehničke metode koje imaju važnu ulogu u maksimizaciji korisničkog angažovanja: kolaborativno filtriranje, učenje potkrepljivanjem i A/B testiranje. One se razlikuju po načinu rada, ali im je zajednička funkcija da platformi omoguće da iz ponašanja korisnika nauči šta zadržava pažnju i da to znanje koristi za dalje prilagođavanje sadržaja.

\subsection{Sistemi preporuke i kolaborativno filtriranje}

Kolaborativno filtriranje jedna je od osnovnih metoda sistema preporuke. Glavna ideja je relativno jednostavna: ako su dva korisnika u prošlosti pokazivala slična interesovanja, postoji verovatnoća da će im i u budućnosti biti zanimljiv sličan sadržaj. Umesto da sistem “razume” sam sadržaj u ljudskom smislu, on posmatra obrasce ponašanja velikog broja korisnika i iz tih obrazaca izvodi preporuke \cite{wen2008collaborative}.

Kolaborativno filtriranje je klasa metoda koje preporučuju stavke korisnicima na osnovu preferencija koje su drugi korisnici izrazili prema tim stavkama. U klasičnom obliku, sistem je zasnovan na matrici u kojoj redovi predstavljaju korisnike, a kolone različite sadržaje — filmove, proizvode, objave, pesme, video-snimke i slično. Vrednosti u matrici predstavljaju neki oblik korisničke reakcije. Ona može biti eksplicitna, kao što je ocena od jedne do pet zvezdica, ali može biti i implicitna, kao što su vreme gledanja, klik, pregled stranice, kupovina ili ponovno vraćanje na isti sadržaj \cite{wen2008collaborative}.

Kod pristupa zasnovanog na korisnicima (engl. \textit{user-based collaborative filtering}), sistem traži korisnike koji su slični posmatranom korisniku, odnosno one koji su gledali iste video-sadržaje, lajkovali slične objave ili kupovali iste proizvode. Na osnovu toga, sistem zaključuje da sadržaj koji su korisnici iz “sličnog klastera” već konzumirali može biti dovoljno relevantan i stimulišući i za posmatranog korisnika. Dakle, logika je: ako su ljudi koji se ponašaju slično kao ja pozitivno reagovali na neki raniji sadržaj, algoritam pretpostavlja da postoji dobra šansa da ću i ja reagovati slično na sadržaj koji planira da mi preporuči.

Osim prethodno pomenutog, postoji i pristup zasnovan na stavkama (engl. \textit{item-based collaborative filtering}). U ovom slučaju korisniku se ne preporučuje sadržaj prvenstveno zato što se svideo drugim korisnicima koji su mu po interakcijama slični, već zato što je taj sadržaj često pozitivno povezan sa sadržajem koji je posmatrani korisnik već konzumirao. Dakle, zaključuje se da su stavke povezane kroz obrasce uparenih pozitivnih reakcija i one ne moraju nužno biti tematski slične. Na primer, dva video-snimka ne moraju govoriti o istoj temi da bi ih algoritam povezao; dovoljno je da veliki broj korisnika često pozitivno reaguje na oba. Ovakav pristup posebno je koristan u kontekstu velikih sistema jer je generalno skalabilniji — lakše je računati sličnosti između stavki sadržaja nego stalno porediti svakog korisnika sa velikim brojem drugih korisnika \cite{wen2008collaborative}.

U kontekstu korisničkog angažovanja, kolaborativno filtriranje je veoma važno jer omogućava platformi da korisniku ne prikazuje samo “generalno popularan sadržaj”, već sadržaj koji je za njega personalizovan i zbog toga verovatno relevantniji. Međutim, važno je razdvojiti korist koju pojedinac može imati od ovakve tehnike od koristi koju od nje imaju same platforme. Upravo tu se otvaraju i društveno problematične implikacije sistema preporuke, uključujući mogućnost stvaranja takozvanog efekta filter-mehura (engl. \textit{filter bubble}), gde personalizacija prestaje da bude samo korisna selekcija sadržaja i postaje mehanizam zatvaranja korisnika u sopstvenu informacionu petlju \cite{nguyen2014filterbubble}.

\subsection{Učenje potkrepljivanjem i dugoročna optimizacija korisničkog angažovanja}

Učenje potkrepljivanjem moderniji je pristup optimizaciji interakcija. Kod kolaborativnog filtriranja sistem pokušava da proceni šta bi korisnik mogao poželeti sada. Kod učenja potkrepljivanjem, sistem preduzima određenu akciju, posmatra reakciju korisnika i na osnovu nje dobija signal nagrade, pa buduće akcije prilagođava novim saznanjima.

U sistemu preporuke, “akcija” može biti izbor sledećeg video-snimka, objave, reklame ili notifikacije. “Nagrada” može biti klik, gledanje do kraja, duže zadržavanje, povratak korisnika na platformu ili neka kombinacija ovih metrika. Ako korisnik reaguje na preporučeni sadržaj, sistem dobija signal da je ta odluka bila uspešna. Ako ne reaguje, sistem dobija slabiji signal i u narednim koracima prilagođava svoje ponašanje.

Tema manipulacije pažnjom i etička pitanja postaju posebno važni kada govorimo o primeni učenja potkrepljivanjem za optimizaciju interakcija, jer sistem u ovom slučaju ne mora da bira sadržaj samo prema trenutnim interesovanjima korisnika. Pošto uči kroz niz preporuka i reakcija, on može početi da uzima u obzir i to kako današnja preporuka utiče na buduće ponašanje korisnika.

U tom kontekstu istraživači uvode pojam manipulacije korisnikom (engl. \textit{user tampering}), kojim opisuju situaciju u kojoj sistem, pokušavajući da dugoročno poveća korisničko angažovanje, ne samo da prati korisnikove postojeće preferencije, već počinje da ih oblikuje. To ne znači da se svaka platforma ponaša na taj način, ali pokazuje zašto učenje potkrepljivanjem može biti posebno osetljivo: ako je “nagrada” vezana za vreme provedeno na platformi, klikove ili učestalost vraćanja, sistem može učiti obrasce koji zadržavaju pažnju i onda kada to nije u najboljem interesu korisnika \cite{kasirzadeh2021user}. Tada optimizacija prestaje da bude neutralna tehnička procedura i prelazi u manipulaciju korisnikom — potencijalno štetnu i po fizičko i po mentalno zdravlje, ponovo zarad profita.

\subsection{A/B testiranje kao eksperimentalna optimizacija platformi}

Pored algoritama za preporuku, digitalne platforme koriste i A/B testiranje kako bi proverile koje verzije proizvoda najbolje ostvaruju zadate ciljeve. A/B testiranje je oblik kontrolisanog eksperimenta u kome se korisnici nasumično dele u grupe: jedna grupa vidi postojeću verziju neke funkcionalnosti, dok druga vidi izmenjenu verziju. Posmatra se koja verzija daje bolje rezultate, mereno na unapred definisanim metrikama \cite{kohavi2017online}.

U kontekstu digitalnih platformi, A/B testiranje može se koristiti za gotovo svaki element korisničkog iskustva:  redosled objava u toku sadržaja (engl. \textit{feed}), raspored dugmadi, boju interfejsa, tekst notifikacije, način prikazivanja komentara, oblik preporuka ili način na koji se korisnik vraća na aplikaciju. Iako neki od prethodno definisanih aspekata deluju trivijalno, kada platforma ima veliki broj korisnika, bilo kakva promena u intenzitetu interakcija može imati značajan efekat.

Ova metoda se koristi zbog toga što konačni sud o promeni nekog aspekta korisničkog iskustva ne prepušta intuiciji dizajnera, programera ili menadžera, već se on donosi na osnovu empirijskih podataka - verzija koja donosi više interakcija biva ocenjena kao uspešnija.

Istraživači naglašavaju da onlajn eksperimenti moraju imati jasno definisanu evaluacionu meru, odnosno kriterijum po kome se procenjuje da li je neka promena zaista bolja, smatrajući da je upravo taj izbor metrike ključan - ako se kao uspeh meri pre svega korisničko angažovanje, onda će platforma sistematski birati verzije koje korisničko angažovanje i povećavaju \cite{kohavi2017online}.

Za temu maksimizacije korisničkog angažovanja posebno je bitno to što A/B testiranje normalizuje stalno eksperimentisanje nad korisničkim iskustvom. Korisnik često nije svestan da je deo eksperimenta, niti zna da se njegova reakcija koristi za odlučivanje o budućem dizajnu platforme. To ne znači da je svako A/B testiranje problematično. Naprotiv, ono može pomoći da se naprave jasniji, brži i korisniji sistemi. Problem ponovo nastaje kada se kao uspeh najviše vrednuje ono što povećava angažovanje, čak i ako to ne znači da je iskustvo za korisnika kvalitetnije.

A/B testiranje je zato tehnički drugačije od kolaborativnog filtriranja i učenja potkrepljivanjem, ali ima sličnu ulogu u ekonomiji pažnje. Ono omogućava platformi da stalno proverava šta korisnike duže zadržava, šta ih vraća i na šta najčešće reaguju. Time se korisničko iskustvo pretvara u prostor neprekidne optimizacije, a korisnik u nepresušni izvor podataka na kome se ta optimizacija proverava.

\subsection{Društvene implikacije tehničke optimizacije}

Kolaborativno filtriranje, učenje potkrepljivanjem i A/B testiranje nisu same po sebi negativne tehnologije. One mogu pomoći korisniku da lakše pronađe relevantan sadržaj, platformi da popravi korisničko iskustvo, a sistemu da se prilagodi realnim potrebama korisnika. Međutim, u poslovnom modelu velikih digitalnih platformi korisničko angažovanje predstavlja jednu od centralnih optimizacionih vrednosti, jer se preko njega posredno maksimizuju pažnja, prikupljanje podataka, oglašavanje, zadržavanje korisnika, a samim tim - prihod. Iako platforme formalno ne optimizuju samo korisničko angažovanje, ono u praksi često postaje ključna posrednička metrika: što korisnik duže ostaje, češće reaguje i redovnije se vraća, to platforma ima više prilika da prikuplja podatke, prikazuje oglase, testira sadržaj i dalje prilagođava sistem \cite{gonzalez2024attention}.

Kod kolaborativnog filtriranja problem nastaje onda kada personalizacija prestane da bude samo pomoć korisniku da se snađe u ogromnoj količini dostupnog sadržaja i postane instrument za što efikasnije zadržavanje njegove pažnje. Sistem koji stalno preporučuje sadržaj sličan onome na šta je korisnik već reagovao može doprineti stvaranju efekta filter-mehura: korisnik sve ređe izlazi iz sopstvene informacione petlje, a sve češće dobija sadržaj koji potvrđuje njegove postojeće reakcije, stavove i slabosti \cite{nguyen2014filterbubble}. U takvom prostoru rasizam, mizoginija, mizandrija, homofobija, ekstremna religijska uverenja, teorije zavere, agresija, mržnja i slični obrasci ne moraju više delovati kao margina interneta. Algoritamski posredovani tok sadržaja može ih pretvoriti u nešto što korisniku izgleda blisko, često, normalno i društveno prihvatljivo.

Kod učenja potkrepljivanjem problem postaje još dublji, jer sistem ne mora samo da procenjuje šta korisnik trenutno želi, već može da uči kako da redosledom preporuka utiče na ono što će korisnik poželeti kasnije. Tada se javlja tzv. manipulacija korisnikom - ne samo praćenje preferencija, već njihovo aktivno oblikovanje. Drugim rečima, sistem može naučiti da manipuliše korisnikom ako mu ta manipulacija pomaže da dugoročno poveća korisničko angažovanje. Sistem učenja potkrepljivanjem može naučiti i da korisnicima rano prikazuje polarizujući sadržaj, kako bi kasnije lakše predviđao i preporučivao sadržaj koji odgovara tako oblikovanim stavovima. Manipulativno ponašanje ne mora nastati uprkos optimizaciji, već upravo kao njena posledica \cite{kasirzadeh2021user}. Ako je nagrada vezana za vreme provedeno na platformi, klikove ili učestalost vraćanja, onda sistem nema razlog da traži ono što je za korisnika dobro, već ono što ga najefikasnije zadržava.

A/B testiranje zatvara ovaj krug. Ako kolaborativno filtriranje bira šta korisnik verovatno želi, a učenje potkrepljivanjem uči kako se njegove buduće reakcije mogu oblikovati, A/B testiranje proverava, u realnom vremenu, koja verzija digitalnog okruženja najbolje radi taj posao. Testira se koja notifikacija ga brže vraća, koji raspored toka sadržaja ga duže zadržava, koja formulacija izaziva više reakcija, koji oblik preporuke probija njegovu pažnju \cite{kohavi2017online}. Korisnik postaje nesvesni učesnik eksperimenata nad sopstvenim ponašanjem: izvor podataka, objekat merenja i pokusni kunić sistema koji proverava šta ga još dublje uvlači.

Zato se zavisnost od digitalnih platformi ne može posmatrati isključivo kao slabost pojedinca ili manjak samokontrole. Naravno da korisnik ima određenu odgovornost za sopstvene navike, ali pre nego što okarakterišemo čoveka kao slabog, potrebno je da ga postavimo u kontekst - sistem koji je tehnički, ekonomski i psihološki podešen da pažnju traži, hvata, zadržava i ponovo aktivira. Idealni konzument takvog sistema nije stabilan, odmoran i potpuno autonoman čovek, već upravo čovek oslabljen teškim životnim situacijama, umorom, usamljenošću, dosadom ili anksioznošću — neko ko se povlači u sadržaj toliko sočno optimizovan da ga uvuče što dublje unutra.

Upravo tu se tehničke metode maksimizacije korisničkog angažovanja spajaju sa logikom ekonomije pažnje i nadzornog kapitalizma. Algoritamska manipulacija pažnjom ne mora da izgleda kao otvorena prinuda, propaganda ili direktna zabrana. Mnogo češće izgleda kao savršeno prirodan tok sadržaja, kao sledeći video koji se sam pokreće, kao notifikacija poslata u pravom trenutku, kao preporuka koja “slučajno” pogađa slabost korisnika. Njena duboka neetičnost nije u jednoj velikoj odluci, već u hiljadama sitnih optimizacija koje pojedinačno deluju neutralno, a zajedno grade sistem u kome je pažnja korisnika sirovina, ponašanje proizvod, a zadržavanje krajnji cilj \cite{gonzalez2024attention}.